\section*{第 \ref{ch:b1:integration} 章 --- 线段上的积分}

\begin{solution}{exo:b1:integration:1}
$\dfrac{1}{x^2-4} = \dfrac{1/4}{x - 2} - \dfrac{1/4}{x+2}$
（遮挡法），于是
\[
\int_0^1 \frac{\dd x}{x^2 - 4}
= \frac14\Bigl[\ln\abs{x-2} - \ln\abs{x+2}\Bigr]_0^1
= \frac14\Bigl(\ln\frac{1}{3} - \ln 1\Bigr) = -\frac{\ln 3}{4}.
\]

用分部积分（$u' = 1$，$v = \ln t$）：$\int_1^{\eu} \ln t\,\dd t =
\bigl[t\ln t\bigr]_1^{\eu} - \int_1^{\eu} \dd t = \eu - (\eu - 1) =
1$。

用分部积分（$u' = \sin t$，$v = t$）：$\int_0^\pi t\sin t\,\dd t =
\bigl[-t\cos t\bigr]_0^\pi + \int_0^\pi \cos t\,\dd t = \pi + 0 =
\pi$。
\end{solution}

\begin{solution}{exo:b1:integration:2}
作代换 $u = t^2 + 1$，$\dd u = 2t\,\dd t$：
\[
\int_0^1 \frac{t\,\dd t}{(t^2+1)^2}
= \frac12 \int_1^2 \frac{\dd u}{u^2}
= \frac12\Bigl[-\frac1u\Bigr]_1^2 = \frac14 .
\]
令 $u = \sin t$：$\int_0^{\pi/2} \cos^3 t\,\dd t = \int_0^{\pi/2}
(1 - \sin^2 t)\cos t\,\dd t = \bigl[\sin t - \frac{\sin^3
t}{3}\bigr]_0^{\pi/2} = 1 - \frac13 = \frac23$。
\end{solution}

\begin{solution}{exo:b1:integration:3}
$u_n = \dfrac1n \sum_{k=1}^{n} \dfrac{1}{1 + (k/n)^2}$ 是
$x \mapsto \frac{1}{1+x^2}$ 在 $\intcc{0}{1}$ 上的黎曼和，于是 $u_n \to
\int_0^1 \frac{\dd x}{1+x^2} = \arctan 1 = \dfrac\pi4$。

$\ln v_n = \dfrac1n \sum_{k=1}^{n} \ln\dfrac{n+k}{n} = \dfrac 1n
\sum_{k=1}^n \ln\Bigl(1 + \dfrac kn\Bigr) \to \int_0^1 \ln(1+x)\,\dd
x = \bigl[(1+x)\ln(1+x) - x\bigr]_0^1 = 2\ln 2 - 1$。于是 $v_n \to
\eu^{2\ln 2 - 1} = \dfrac 4\eu$。（检验此恒等式：
$\frac{(2n)!}{n!\,n^n} = \prod_{k=1}^{n} \frac{n+k}{n}$。）
\end{solution}

\begin{solution}{exo:b1:integration:4}
极限为 $0$。设 $M = \sup \abs f$ 且 $\varepsilon \in
\intoo{0}{1}$。在 $1 - \varepsilon$ 处分段：
\[
\Bigl| \int_0^1 x^n f \Bigr|
\leq \int_0^{1 - \varepsilon} x^n \abs f + \int_{1-\varepsilon}^1
x^n \abs f
\leq M\,(1-\varepsilon)^n + M\varepsilon .
\]
由于 $(1 - \varepsilon)^n \to 0$（见\cref{exo:b1:seq:3}），左边的上极限
$\leq M\varepsilon$ 对每个 $\varepsilon$ 都成立：
故积分趋于 $0$。
\end{solution}

\begin{solution}{exo:b1:integration:5}
$Q(\lambda) = \int_a^b (f + \lambda g)^2 = \int f^2 + 2\lambda \int
fg + \lambda^2 \int g^2 \geq 0$ 对所有 $\lambda$ 成立。若 $\int g^2 =
0$，则 $g = 0$（严格正性，
\cref{thm:b1:integration:props}~(4)），于是不等式为 $0 \leq
0$。否则 $Q$ 是真正的二次多项式，处处 $\geq 0$，其
判别式 $\leq 0$，即 $\bigl(\int fg\bigr)^2 \leq \int
f^2 \int g^2$。

等号成立当且仅当判别式为零，亦即 $Q(\lambda_0) = 0$ 对某个
$\lambda_0$，即 $\int (f + \lambda_0 g)^2 = 0$，亦即（再由严格正性
）$f = -\lambda_0 g$，故
等号成立当且仅当 $f$ 与 $g$ 成比例。
\end{solution}

\begin{solution}{exo:b1:integration:6}
设 $\Phi(a) = \int_a^{a+T} f$。由微积分基本定理
（\cref{thm:b1:integration:ftc}），$\Phi$ 可导，且
$\Phi'(a) = f(a + T) - f(a) = 0$，故为常数。

对 $x > 0$，记 $x = nT + r$，$0 \leq r < T$（$n = \lfloor x/T
\rfloor$）。由沙勒关系：
\[
\int_0^x f = n \int_0^T f + \int_{nT}^{nT + r} f,
\qquad
\Bigl| \int_{nT}^{nT+r} f \Bigr| \leq T \sup_{\intcc{0}{T}} \abs f
= C .
\]
于是 $\frac 1x \int_0^x f = \frac{nT}{x}\cdot\frac 1T \int_0^T f +
O\bigl(\frac 1x\bigr)$，且 $\frac{nT}{x} \to 1$，极限为
$\frac1T \int_0^T f$。
\end{solution}

\begin{solution}{exo:b1:integration:7}
设 $g(x) = f(x) - x$，它连续，且
\[
\int_0^1 g = \int_0^1 f - \frac12 = 0 .
\]
若 $g$ 处处不为零，则由介值定理，
必保持定号（区间上同时取两种符号的连续函数必有零点）；
不妨设 $g > 0$。再由严格正性
（\cref{thm:b1:integration:props}~(4)，对 $g > 0$ 应用，得
$\int g > 0$），与 $\int g = 0$ 矛盾。故 $g(c) = 0$ 对某个
$c$，即 $f(c) = c$。
\end{solution}

\begin{solution}{exo:b1:integration:8}
\begin{enumerate}
  \item 用分部积分，取 $u' = \sin t$，$v = \sin^{n-1} t$：
        \[
        W_n = \bigl[-\cos t\sin^{n-1}t\bigr]_0^{\pi/2}
        + (n-1)\int_0^{\pi/2} \cos^2 t\,\sin^{n-2} t\,\dd t
        = (n-1)(W_{n-2} - W_n),
        \]
        于是 $nW_n = (n-1)W_{n-2}$。由 $W_0 = \frac\pi2$，$W_1 = 1$：
        \[
        W_{2p} = \frac{(2p-1)(2p-3)\cdots 1}{(2p)(2p-2)\cdots 2}\,
        \frac{\pi}{2} = \frac{(2p)!}{4^p (p!)^2}\,\frac\pi2,
        \qquad
        W_{2p+1} = \frac{(2p)(2p-2)\cdots 2}{(2p+1)(2p-1)\cdots 3}
        = \frac{4^p (p!)^2}{(2p+1)!} .
        \]
  \item 在 $\intoo{0}{\frac\pi2}$ 上，$0 < \sin t < 1$，故 $\sin^{n+1} <
        \sin^n$，$(W_n)$ 严格递减且为正。
        由递推关系作夹逼：
        \[
        \frac{n}{n+1} = \frac{W_{n+1}}{W_{n-1}}
        \leq \frac{W_{n+1}}{W_n} \leq 1
        \quad\implies\quad \frac{W_{n+1}}{W_n} \to 1 .
        \]
        不变量 $a_n = (n+1)W_{n+1}W_n$ 满足 $a_n = a_{n-1}$，
        由递推关系 $(n+1)W_{n+1} = nW_{n-1}$，故 $a_n = a_0 =
        1 \cdot W_1 W_0 = \frac\pi2$。于是
        \[
        n W_n^2 \sim (n+1) W_{n+1} W_n = \frac\pi2
        \quad\implies\quad
        W_n \sim \sqrt{\frac{\pi}{2n}} .
        \]
\end{enumerate}
\end{solution}

\begin{solution}{exo:b1:integration:9}
\begin{enumerate}
  \item 在 $\intoo{0}{\pi}$ 上：$x > 0$，$a - bx = b(\frac ab - x) =
        b(\pi - x) > 0$，$\sin x > 0$，故被积函数 $> 0$，
        从而 $I_n > 0$（严格正性）。估计：在
        $\intcc{0}{\pi}$ 上，$x \leq \pi$，$a - bx \leq a$，故 $P
        \leq \frac{\pi^n a^n}{n!}$，且 $I_n \leq
        \pi\,\frac{(\pi a)^n}{n!}$，它趋于 $0$（阶乘
        增长快于几何项：它是
        收敛指数级数的通项，参见
        \cref{ex:b1:seq:e}）；特别地，$I_n < 1$ 对充分大的 $n$ 成立。
  \item 展开 $x^n(a - bx)^n = \sum_{j=0}^{n} \binom nj a^{n-j}
        (-b)^j x^{n+j}$，故 $P = \frac{1}{n!}\sum_j c_j x^{n+j}$，
        其中 $c_j$ 为整数。于是 $P^{(k)}(0) = 0$，当 $k < n$ 时
        （由零点阶数），且对 $n \leq k \leq 2n$，$P^{(k)}(0) =
        \frac{k!}{n!} c_{k-n}$，它是整数，因为 $n! \mid k!$。
        此外 $P(\pi - x) = P(x)$（代入：$\pi - x$ 交换两个因子，
        利用 $a - b(\pi - x) = bx$），故 $P^{(k)}(\pi)
        = \pm P^{(k)}(0)$，同样为整数。
  \item 令 $Q = P - P'' + P^{(4)} - \dots$（有限和，$P$ 的次数为
        $2n$）：$Q + Q'' = P$，且
        \[
        \bigl(Q'\sin x - Q\cos x\bigr)' = (Q + Q'')\sin x = P\sin x .
        \]
        于是 $I_n = \bigl[Q'\sin x - Q\cos x\bigr]_0^{\pi} = Q(\pi)
        + Q(0)$，它是若干 $P^{(2k)}$ 在 $0$ 与 $\pi$ 处的值之和，
        由 (2) 知其为整数。
  \item 当 $n$ 充分大时，$I_n$ 为整数，且 $0 < I_n < 1$：
        矛盾。假设 $\pi = \frac ab$ 不成立，故 $\pi$ 是
        无理数。
\end{enumerate}
\end{solution}

\begin{solution}{exo:b1:integration:10}
$C^1$ 情形：分部积分，
\[
\int_a^b f(t)\sin\lambda t\,\dd t
= \Bigl[-f(t)\frac{\cos\lambda t}{\lambda}\Bigr]_a^b
+ \frac{1}{\lambda}\int_a^b f'(t)\cos\lambda t\,\dd t,
\]
其绝对值不超过 $\frac{2\sup\abs f + (b - a)\sup\abs{f'}}
{\lambda} \to 0$.

连续情形：取 $\varepsilon > 0$，并取一个阶梯函数
$\varphi$，使得 $\abs{f - \varphi} \leq \varepsilon$
（由\cref{thm:b1:integration:approx}可知存在 $\varphi \leq f \leq
\psi$，且间隙 $\leq\varepsilon$；取 $\varphi$）。于是
\[
\Bigl|\int_a^b f\sin\lambda t\Bigr|
\leq \int_a^b \abs{f - \varphi}
+ \Bigl|\int_a^b \varphi \sin\lambda t\Bigr|
\leq (b-a)\varepsilon
+ \sum_i \abs{c_i}\,\Bigl|\int_{x_{i-1}}^{x_i} \sin\lambda t\,\dd
t\Bigr| ,
\]
且每个 $\bigl|\int \sin \lambda t\bigr| = \bigl|\frac{\cos\lambda
x_{i-1} - \cos\lambda x_i}{\lambda}\bigr| \leq \frac{2}{\lambda}$:
第二项趋于 $0$。于是上极限 $\leq
(b-a)\varepsilon$（对每个 $\varepsilon$ 成立），故极限为 $0$。
\end{solution}

\begin{solution}{exo:b1:integration:11}
设 $m = \min f$ 与 $M = \max f$，它们由
最值定理取到。由于 $g \geq 0$，有 $m\,g \leq fg \leq M\,g$，于是由
单调性得
\[
m \int_a^b g \;\leq\; \int_a^b fg \;\leq\; M \int_a^b g .
\]
若 $\int_a^b g = 0$，则由严格正性
（\cref{thm:b1:integration:props}~(4)）推出 $g \equiv 0$，两侧
均为零，任意 $c$ 皆可。否则 $t = \frac{\int
fg}{\int g}$ 落在 $\intcc{m}{M} = f(\intcc{a}{b})$
由\cref{thm:b1:continuity:evt,thm:b1:continuity:ivt}可知，故 $t =
f(c)$，即存在某个 $c$。

符号至关重要：在 $\intcc{-1}{1}$ 上取 $f(t) = g(t) = t$，则 $\int
fg = \int_{-1}^1 t^2 = \frac23$，而 $f(c)\int_{-1}^1 t\,\dd t
= 0$，对每个 $c$ 都成立。
\end{solution}

\begin{solution}{exo:b1:integration:12}
若 $f \equiv 0$，则结论平凡成立（每一点都是零点）。故
设 $f \not\equiv 0$，并假设它至多有 $n$ 个不同
零点，位于 $\intoo{a}{b}$ 内；记 $z_1 < \dots < z_m$（$m \leq n$）为
这些使 $f$ \emph{变号}的零点（可能一个也没有）。令
$P(t) = \prod_{i=1}^{m}(t - z_i)$（空积 $= 1$），其次数为
$m \leq n$。在由各 $z_i$ 分出的每个子区间上，
$f$ 与 $P$ 均保持定号，且跨过
某个 $z_i$ 时两者都变号：乘积 $fP$ 在整个
$\intoo{a}{b}$ 上定号。由于连续、不恒为零，
故 $\bigl|\int_a^b fP\bigr| > 0$（严格
正性用于 $\abs{fP}$）。但 $\int f P$ 是
关于矩 $\int f\,t^k$（$k \leq n$）的线性组合，而这些矩全为零：
矛盾。因此 $f$ 至少有 $n + 1$ 个不同零点，位于
$\intoo{a}{b}$ 内。
\end{solution}

\begin{solution}{pb:b1:integration:1}
\textbf{1.} 设 $N = \lceil 2c \rceil$。当 $n \geq N$ 时：
$\frac{c^{n+1}/(n+1)!}{c^n/n!} = \frac{c}{n+1} \leq \frac12$，故
$0 < \frac{c^n}{n!} \leq \frac{c^N}{N!}\,2^{-(n - N)} \to 0$:
夹逼。

\textbf{2.} 固定 $k$，对 $m$ 作归纳。当 $m = 0$ 时：$\int_0^1
x^k = \frac{1}{k+1} = \frac{k!\,0!}{(k+1)!}$。归纳步，分部积分
($u = (1-x)^m$, $v' = x^k$):
\[
\int_0^1 x^k (1-x)^m \dd x
= \Bigl[\frac{x^{k+1}}{k+1}(1-x)^m\Bigr]_0^1
+ \frac{m}{k+1}\int_0^1 x^{k+1}(1-x)^{m-1}\dd x
= \frac{m}{k+1}\cdot\frac{(k+1)!\,(m-1)!}{(k+m+1)!} ,
\]
即 $\frac{k!\,m!}{(k+m+1)!}$。

\textbf{3.} $k = m = n$: $\int_0^1 (x(1-x))^n =
\frac{(n!)^2}{(2n+1)!} = \frac{1}{(2n+1)\binom{2n}{n}}$。由于
$x(1-x) \leq \frac14$ 在 $\intcc{0}{1}$ 上成立，因此积分 $\leq
4^{-n}$，从而 $\binom{2n}{n} \geq \frac{4^n}{2n+1}$ ---
这与 $\binom{2n}{n}^{1/n} \to 4$ 相符
（见\cref{pb:b1:seq:1}）。

\textbf{4.} 被积函数连续，$\geq 0$，且在
$\intoo{0}{1}$ 上为正，因而不恒为零，故其积分
$> 0$（由\cref{thm:b1:integration:props}~(4)）。上界估计：
$(x(1-x))^n \leq 4^{-n}$，$g \leq \sup\abs g$，再由
单调性即得。

\textbf{5.} $A_0 = \eu - 1$；$A_1 = [x\eu^x]_0^1 - \int_0^1
\eu^x = \eu - (\eu - 1) = 1$。由分部积分，$A_n = [x^n \eu^x]_0^1 -
n\int_0^1 x^{n-1}\eu^x = \eu - n A_{n-1}$。估计：被积函数
为正，故 $A_n > 0$；又由 $\eu^x \leq \eu$ 得 $A_n \leq
\eu\int_0^1 x^n = \frac{\eu}{n+1}$。

\textbf{6.} $A_0 = -1 + 1\cdot\eu$。若 $A_{n-1} = \alpha_{n-1} +
\beta_{n-1}\eu$，其系数皆为整数，则
\[
A_n = \eu - n\alpha_{n-1} - n\beta_{n-1}\eu
= \underbrace{(-n\,\alpha_{n-1})}_{\alpha_n}
+ \underbrace{(1 - n\,\beta_{n-1})}_{\beta_n}\,\eu ,
\]
二者皆为整数。

\textbf{7.} 若 $\eu = \frac pq$，则 $q A_n = q\alpha_n + p\beta_n
\in \Z$，且 $0 < qA_n \leq \frac{q\eu}{n+1} < 1$ 对充分大的 $n$ 成立：
一个严格介于 $0$ 与 $1$ 之间的整数，不可能 --- 故 $\eu
\notin \Q$。在\cref{exo:b1:seq:9}中，被夹住的整数是
$q!\,\frac pq - q!\,a_q$；此处为 $qA_n$：同一陷阱，
燃料来自积分。

\textbf{8.} 取 $n = 0$，则 $R_0 = \int_0^1 \eu^t = \eu - 1$，故 $\eu =
1 + R_0$。分部积分（取 $u = \eu^t$，$v = -\frac{(1-t)^{n+1}}{n+1}$）：
\[
R_n = \frac{1}{n!}\Bigl(\Bigl[-\frac{(1-t)^{n+1}}{n+1}
\eu^t\Bigr]_0^1 + \frac{1}{n+1}\int_0^1 (1-t)^{n+1}\eu^t\Bigr)
= \frac{1}{(n+1)!} + R_{n+1} ,
\]
于是该公式由 $n$ 递推到 $n + 1$。估计：$1 \leq
\eu^t \leq \eu$ 在 $\intcc{0}{1}$ 上成立，且 $\int_0^1 (1-t)^n =
\frac{1}{n+1}$，于是 $\frac{1}{(n+1)!} \leq R_n \leq
\frac{\eu}{(n+1)!}$。

\textbf{9.} $\sum_{k=0}^{9} \frac{1}{k!} = \frac{986410}{362880}
= 2.71828152\dots$，且
\[
\frac{1}{10!} = 2.76\cdot10^{-7} \leq R_9 \leq
\frac{2.75}{10!} = 7.58\cdot10^{-7} ,
\]
故 $2.7182818 \leq \eu \leq 2.7182823$；由此严格地得 $\eu =
2.718282$（精确到六位小数；真值为 $2.7182818\dots$）。

\textbf{10.} $f(1 - x) = \frac{(1-x)^n x^n}{n!} = f(x)$。在
$\intoo{0}{1}$ 上 $0 < x(1-x) \leq \frac14$，故 $0 < f \leq
\frac{1}{4^n\,n!}$。

\textbf{11.} $x^n(1-x)^n = \sum_{j=0}^{n} (-1)^j\binom
nj\,x^{n+j}$，故 $f = \frac{1}{n!}\sum_j c_j\,x^{n+j}$，其中
$c_j \in \Z$。因此 $f^{(k)}(0) = 0$，当 $k < n$ 或 $k > 2n$ 时，
当 $n \leq k \leq 2n$ 时：$f^{(k)}(0) = \frac{k!}{n!}\,
c_{k-n}$，它是整数，因为 $n! \mid k!$。由对称性得
$f^{(k)}(1) = (-1)^k f^{(k)}(0) \in \Z$。

\textbf{12.} $b^n \pi^{2n-2k} = b^n\bigl(\frac
ab\bigr)^{\,n-k} = a^{\,n-k}\,b^{\,k} \in \Z$，故 $G(0) =
\sum_k (-1)^k a^{n-k} b^k f^{(2k)}(0)$，同理 $G(1)$ 亦为
整数，由第 11 题可知。

\textbf{13.} 在 $\pi^2 G + G''$ 中，第 $k$ 项（取自 $\pi^2 G$）
携带 $\pi^{2n-2k+2}f^{(2k)}$，而第 $j = k - 1$ 项
（属于 $G''$）携带 $(-1)^{k-1}\pi^{2n-2k+2}f^{(2k)}$：各项全部相消，
只剩第一和式中 $k = 0$ 与第二和式中 $j = n$ 的项，
即
\[
G'' + \pi^2 G = b^n\bigl(\pi^{2n+2} f + (-1)^n f^{(2n+2)}\bigr)
= b^n \pi^{2n+2} f = \pi^2 a^n f
\]
（$f$ 的次数为 $2n$，故 $f^{(2n+2)} = 0$；又 $b^n\pi^{2n} =
a^n$）。于是
\[
\bigl(G'\sin\pi x - \pi G\cos\pi x\bigr)'
= (G'' + \pi^2 G)\sin \pi x = \pi^2 a^n f\sin\pi x .
\]

\textbf{14.} 在 $\intcc{0}{1}$ 上积分：
\[
\pi^2 a^n \int_0^1 f\sin(\pi x)\,\dd x
= \bigl[G'\sin\pi x - \pi G\cos\pi x\bigr]_0^1
= \pi\bigl(G(1) + G(0)\bigr) ,
\]
故 $\pi a^n \int_0^1 f\sin\pi x = G(0) + G(1) \in \Z$。在
$\intoo{0}{1}$ 上，$f > 0$ 且 $\sin\pi x > 0$，因此左边
为正，故 $G(0) + G(1) \geq 1$；又 $\sin \leq 1$，由第 10 题，
可将其上界估为 $\frac{\pi a^n}{4^n\,n!}$。

\textbf{15.} 由第 1 题（取 $c = \frac a4$），
$\frac{\pi a^n}{4^n n!} \to 0$：对充分大的 $n$，它 $< 1$，
与 $G(0) + G(1) \geq 1$ 矛盾。故不存在分数 $\frac ab$
与 $\pi^2$ 相等，此即勒让德定理。若 $\pi$ 为有理数，
则 $\pi^2$ 亦为有理数，故 $\pi \notin \Q$ --- 且该蕴含
只能单向成立（$\sqrt2$ 为无理数，其平方却为有理数），
这正是 $\pi^2 \notin \Q$ 严格强于
\cref{exo:b1:integration:9} 的原因。

\textbf{16.} 整值性：第 11--12 题（端点
导数）；微小性：第 10 题与第 1 题；桥梁：第
13--14 题（逐项相消的两次分部积分）。若去掉
$\frac{1}{n!}$，端点数据仍为整数（更易看出），但界变为
$\frac{\pi a^n}{4^n}$，它趋于 $0$ 仅当
$a < 4$ 时成立 --- 而每个候选者都有 $a =
b\,\pi^2 > 9$。阶乘正是跑赢
几何增长 $a^n$：没有阶乘，就没有定理。

\textbf{17.} 当 $a \leq 10$ 时，整数 $G(0) + G(1)$
为正，且至多为 $\pi\,(10/4)^n/n!$。当 $n = 7$ 时，
$2.5^7 = 610.35\dots$，故界为 $\frac{\pi \times
610.35}{5040} \approx 0.38 < 1$（当 $n = 6$ 时仍为
$1.07$）：矛盾在第七步
明确显现。

\textbf{18.} 两次分部积分：
\[
\int_0^1 x(1-x)\sin\pi x\,\dd x
= \frac1\pi\int_0^1 (1 - 2x)\cos\pi x\,\dd x
= \frac{2}{\pi^2}\int_0^1 \sin\pi x\,\dd x
= \frac{2}{\pi^2}\cdot\frac{2}{\pi} = \frac{4}{\pi^3}
\]
（边界项为零：$x(1-x)$ 在 $0, 1$ 处为零，$\sin\pi x$
在 $0, 1$ 处亦为零）。形式上，$n = 1$ 的裂项相消
（不假设 $\pi$ 的取值）写作 $\pi^3\int_0^1 f_1\sin\pi x = -(f_1''(0) +
f_1''(1))$，其中 $f_1 = x(1 - x)$，$f_1'' = -2$，右端为 $4$
--- 两种计算结果一致。

\textbf{19.} $\eu^x$ 满足 $y' = y$，$\sin\pi x$ 满足
$y'' = -\pi^2 y$：二者都是常系数线性方程，
因此分部积分会使核循环回到自身，并
使全部边界数据仍落在 $\Z + \Z\eu$ 中（相应地，整系数
$\pi^2$ 的多项式）。这一封闭性正是
机器所需要的。借助同时适配于若干点的核加以细化，
同一机制便得出埃尔米特定理（$\eu$
为超越数，1873 年）与林德曼定理（$\pi$ 为超越数，
1882 年） --- 超出本册范围。

\textbf{20.} 作多项式除法（或乘回验算）：
\[
x^4(1-x)^4 = (x^6 - 4x^5 + 5x^4 - 4x^2 + 4)(1 + x^2) - 4 .
\]
对所示恒等式两边除以 $1 + x^2$ 后积分：
\[
\int_0^1 \frac{x^4(1-x)^4}{1+x^2}\dd x
= \Bigl(\frac17 - \frac46 + 1 - \frac43 + 4\Bigr)
- 4\arctan 1
= \frac{22}{7} - \pi ,
\]
利用 $\frac17 - \frac23 + 1 - \frac43 + 4 = 3 + \frac17$ 及
$\arctan 1 = \frac\pi4$。

\textbf{21.} 被积函数在
$\intoo{0}{1}$ 上连续且为正，故积分 $> 0$，于是 $\pi < \frac{22}{7}$。
此外，$\frac12 \leq \frac{1}{1+x^2} \leq 1$ 在
$\intcc{0}{1}$ 上成立，且 $\int_0^1 (x(1-x))^4 = \frac{(4!)^2}{9!} =
\frac{1}{630}$（第 3 题）：
\[
\frac{1}{1260} \leq \frac{22}{7} - \pi \leq \frac{1}{630}
\quad\Longrightarrow\quad
3.14126 < \frac{22}{7} - \frac{1}{630} \leq \pi \leq
\frac{22}{7} - \frac{1}{1260} < 3.14207 .
\]

\textbf{22.} 模 $x^2 + 1$：$x^2 \equiv -1$，故 $x^{4m} =
(x^2)^{2m} \equiv 1$，且 $(1 - x)^2 = 1 - 2x + x^2 \equiv -2x$，
于是 $(1-x)^{4m} \equiv (-2x)^{2m} = 4^m (x^2)^m \equiv
(-4)^m$：余式为常数 $(-4)^m$，而
商 $Q_m$ 的系数为整数（对首一的
整系数多项式作除法）。将恒等式除以 $1 + x^2$，
并积分：
\[
J_m := \int_0^1 \frac{(x(1-x))^{4m}}{1+x^2}\dd x
= s_m + (-4)^m\,\frac{\pi}{4},
\qquad s_m = \int_0^1 Q_m \in \Q .
\]
解出 $\pi$：令 $r_m = (-1)^{m+1}\,4^{\,1-m} s_m \in
\Q$, $\;\abs{\pi - r_m} = 4^{\,1-m} J_m \leq
4^{\,1-m}\cdot4^{-4m} = 4^{\,1-5m}$。当 $m = 1$ 时，$s_1 =
\frac{22}{7}$，$r_1 = \frac{22}{7}$，界为 $4^{-4} =
\frac{1}{256}$ --- 与第 20--21 题的结果一致。

\textbf{23.} $\bigl|\pi - \frac{22}{7}\bigr| = \frac{22}{7} -
\pi \approx 1.26\cdot10^{-3}$，较
$2$ 阶基准 $\frac{1}{7^2} \approx 2.0\cdot10^{-2}$ 好 16 倍，
由狄利克雷所保证（见\cref{pb:b1:derivative:1}，第 4 题）；
而 $\bigl|\pi - \frac{355}{113}\bigr| \approx 2.7\cdot10^{-7}$
优于 $\frac{1}{113^2} \approx 7.8\cdot10^{-5}$，优势
约为 $\approx 300$ 倍。已证结论中并无矛盾：刘维尔不等式
只给出逼近误差的\emph{下界}，且仅对
代数数成立，而在本阶段对 $\pi$ 并无此类界可用
--- $\pi$ 的逼近可以好得惊人，
这并无禁止。

\textbf{24.} 引理：设 $x = \frac pq$，且 $0 < \abs{a_n +
b_n x} \to 0$。则 $\abs{a_n + b_n x} = \frac{\abs{q a_n + p
b_n}}{q}$，其中 $q a_n + p b_n$ 为非零整数（非零是因为
绝对值 $> 0$）。故 $\abs{a_n + b_n x}
\geq \frac1q$ 对每个 $n$ 成立，与收敛于
$0$ 相矛盾。实例：第 7 题（$x = \eu$，$a_n = \alpha_n$，$b_n =
\beta_n$）；\cref{exo:b1:seq:9}（仍取 $x = \eu$，其中 $a_q =
-q!\sum_{k \leq q}\frac{1}{k!}$，$b_q = q!$）；以及
\cref{pb:b1:derivative:1}，第 1 题，即其几何形式。
在第 III 部分及\cref{exo:b1:integration:9}中，陷阱
恰好发生于矛盾的\emph{内部}：假设其为有理数，便把某个表达式
化为整数，而分析随即将其压入
$\intoo{0}{1}$ 之中 --- 原理相同，只是形式不同。

\textbf{25.} (i) 整性体现于多项式的端点演算中，
即第 6、11--12 题；小性存在于
$\sup$ 单调性所给出的界中（第 4、10 题），桥梁存在于
分部积分（第 8、13--14 题） --- 三者皆
为本章定理。(ii) 积分提供了中值定理所不能给出的
东西：解析对象与算术数据之间的一个\emph{精确}恒等式
（是等式，而不只是
含未知数 $c$ 的不等式），这正是机器
能够达到 $\pi^2$，而\cref{pb:b1:derivative:1}
只得到逼近指数。(iii) 所提取的结论：$\eu \notin \Q$，
$\pi^2 \notin \Q$（从而 $\pi \notin \Q$），$\eu = 2.718282$
经证实，$\frac{22}{7} - \frac{1}{630} \leq \pi \leq
\frac{22}{7} - \frac{1}{1260}$，且 $\binom{2n}{n} \geq
\frac{4^n}{2n+1}$。(iv) 前沿：埃尔米特与林德曼将
同一机器推向超越性；阿佩里（1979 年）将
整数陷阱用于 $\zeta(3)$ --- 机器仍在产出
二十世纪的数学。
\end{solution}
